%% =====================================================================
%%  B4-T-19-02 -- what B4 contributes to "Unequal Exchange"
%%  -------------------------------------------------------------------
%%  LAYER A: APPEND-ONLY. Written once, then left alone. Material found
%%  only here sits in the `unique` slot and is protected by rule R10.
%% =====================================================================
%% @kind: source
%% @id: B4-T-19-02
%% @book: B4
%% @topic: T-19-02
%% @locator: ch. 2, pp. 43--86
%% @status: distilled
%% @unique: yes
%% @integrated: yes
%% @updated: 2026-09-19

\dossier{B4}{T-19-02}{\bookshort{B4}}{ch. 2, pp. 43--86}

\begin{distilled}
  \item An obligation set up by a small benefit can be discharged by a substantially larger one; the rule demands reciprocity, not equivalence.
  \item The uncomfortable state of owing is itself a motive to comply, since agreeing ends it faster than negotiating does.
  \item The sample, the small gift and the minor courtesy are all described as opening a position from which a much larger request is later made.
\end{distilled}

\begin{unique}
  \item The absence of an exchange rate as the mechanism, and the consequent conclusion that the creditor nominates the size of the repayment.
  \item Discomfort as an independent driver of compliance, separate from the social marking.
\end{unique}
